% Source: physical PDF pp. 145--150 (printed pp. 122--127), from the
% Section 27.2 heading through the paragraph immediately before Section 27.3.
% Inventory: Eqs. (27.2.1)--(27.2.20), seven unnumbered display groups, one
% linked citation, one source footnote, and no figures, tables, or dividers.
% Uncertainties: none.

\providecommand{\SuperD}{\mathcal{D}}

\subsection{Gauge-Invariant Action for Abelian Gauge Superfields}

We must now consider how to construct a gauge-invariant supersymmetric
action for the gauge superfields \(V^A(x,\theta)\) that contain the gauge
fields \(V_\mu^A(x)\). In order to motivate this construction we will first
consider the case of a single Abelian gauge field (dropping the superscript
\(A\)), and then return to the general case in the next section.

In an Abelian gauge theory like quantum electrodynamics, the gauge-invariant
field constructed from \(V_\mu(x)\) is the familiar field-strength tensor
\begin{equation}
  f_{\mu\nu}(x)
  =\partial_\mu V_\nu(x)-\partial_\nu V_\mu(x).
  \tag{27.2.1}\label{eq:27.2.1}
\end{equation}
The supersymmetry transformation rule for \(f_{\mu\nu}(x)\) is then given
by the transformation rule \eqref{eq:26.2.15} for \(V_\mu(x)\) as
\begin{equation}
  \delta f_{\mu\nu}
  =-\ii\left[
    \alpha\sigma_\nu\partial_\mu\bar\lambda
    +
    \bar\alpha\bar\sigma_\nu\partial_\mu\lambda
    -(\mu\leftrightarrow\nu)
  \right].
  \tag{27.2.2}\label{eq:27.2.2}
\end{equation}
Eq.~\eqref{eq:26.2.16} gives the transformation rule for \(\lambda(x)\) as
\begin{equation}
  \delta\lambda_\alpha
  =f_{\mu\nu}(\sigma^{\mu\nu}\alpha)_\alpha
    +\ii D\alpha_\alpha,
  \tag{27.2.3}\label{eq:27.2.3}
\end{equation}
while Eq.~\eqref{eq:26.2.17} gives the transformation rule for \(D(x)\):
\begin{equation}
  \delta D
  =\alpha\sigma^\mu\partial_\mu\bar\lambda
   -\bar\alpha\bar\sigma^\mu\partial_\mu\lambda.
  \tag{27.2.4}\label{eq:27.2.4}
\end{equation}
None of this depends on whether or not the superfield \(V^A(x,\theta)\) is
taken to be in Wess--Zumino gauge. We see that the fields
\(f_{\mu\nu}(x)\), \(\lambda(x)\), and \(D(x)\) form a complete
supersymmetry multiplet.

It is not hard to construct a suitable kinematic Lagrangian density for
the fields of this supermultiplet. The only Lorentz-invariant, parity
conserving, and gauge-invariant functions of these fields with
dimensionality four are \(f_{\mu\nu}f^{\mu\nu}\),
\(\ii\bar\lambda\bar\sigma^\mu\partial_\mu\lambda\), and \(D^2\). We can
make \(V^\mu\) a
conventionally normalized vector field by taking the coefficient of
\(f_{\mu\nu}f^{\mu\nu}\) to be \(-1/4\), so we may tentatively take the
kinematic Lagrangian density as
\[
  \mathcal{L}_{\mathrm{gauge}}
  =-\frac14 f_{\mu\nu}f^{\mu\nu}
  +2ic_\lambda\bar\lambda\bar\sigma^\mu\partial_\mu\lambda
  -c_DD^2,
\]
with coefficients \(c_\lambda\) and \(c_D\) to be determined from the
condition that \(\int\mathcal{L}_{\mathrm{gauge}}\dd^4x\) is supersymmetric.
Using Eqs.~\eqref{eq:27.2.2}--\eqref{eq:27.2.4}, an infinitesimal
supersymmetry transformation changes the operators in the Lagrangian by
\[
  \begin{aligned}
  \delta\bigl(f_{\mu\nu}f^{\mu\nu}\bigr)
  &=-2if^{\mu\nu}\left[
    \alpha(\sigma_\nu\partial_\mu-\sigma_\mu\partial_\nu)\bar\lambda
    +\bar\alpha(\bar\sigma_\nu\partial_\mu
      -\bar\sigma_\mu\partial_\nu)\lambda\right],
  \\
  \delta\!\left(2\ii\bar\lambda\bar\sigma^\rho
      \partial_\rho\lambda\right)
  &=-2\left[
    \alpha\bigl(f_{\mu\nu}\sigma^{\mu\nu}+\ii D\bigr)
      \sigma^\rho\partial_\rho\bar\lambda
    +\bar\alpha\bigl(f_{\mu\nu}\bar\sigma^{\mu\nu}-\ii D\bigr)
      \bar\sigma^\rho\partial_\rho\lambda
  \right],
  \\
  \delta D^2&=2D\left(
    \alpha\sigma^\mu\partial_\mu\bar\lambda
    -\bar\alpha\bar\sigma^\mu\partial_\mu\lambda\right),
  \end{aligned}
\]
in which we drop derivative terms that do not contribute to the variation
of the action. In order to see how these terms cancel, it is necessary to
use the sigma-matrix identities\footnote{They follow immediately by
expanding the products of Pauli matrices; the signs of the
Levi--Civita terms are also fixed by the definitions in the Notation
section.}
\begin{equation}
  \begin{aligned}
  \sigma^{\mu\nu}\sigma^\rho
  &=\frac12\left(
    \eta^{\mu\rho}\sigma^\nu-\eta^{\nu\rho}\sigma^\mu
    -\ii\epsilon^{\mu\nu\rho\sigma}\sigma_\sigma\right),\\
  \bar\sigma^{\mu\nu}\bar\sigma^\rho
  &=\frac12\left(
    \eta^{\mu\rho}\bar\sigma^\nu-\eta^{\nu\rho}\bar\sigma^\mu
    +\ii\epsilon^{\mu\nu\rho\sigma}\bar\sigma_\sigma\right).
  \end{aligned}
  \tag{27.2.5}\label{eq:27.2.5}
\end{equation}
The terms proportional to
\(\epsilon^{\mu\nu\rho\sigma}f_{\mu\nu}
\alpha\sigma_\sigma\partial_\rho\bar\lambda\) and its conjugate
does not contribute to \(\int \dd^4x\,\delta\mathcal{L}\), because
integrating by parts yields a contribution proportional to
\(\epsilon^{\mu\nu\rho\sigma}\partial_\rho f_{\mu\nu}\), which vanishes
for an \(f_{\mu\nu}\) of the form \eqref{eq:27.2.1}. This identity then
allows us to rewrite the variation of the \(\lambda\)-term as
\[
  \begin{aligned}
  \delta\!\left(2\ii\bar\lambda\bar\sigma^\rho
      \partial_\rho\lambda\right)
  ={}&-if^{\mu\nu}\left[
    \alpha(\sigma_\nu\partial_\mu-\sigma_\mu\partial_\nu)\bar\lambda
    +\bar\alpha(\bar\sigma_\nu\partial_\mu
      -\bar\sigma_\mu\partial_\nu)\lambda\right]
  \\
  &-2D\left(
    \alpha\sigma^\mu\partial_\mu\bar\lambda
    -\bar\alpha\bar\sigma^\mu\partial_\mu\lambda\right).
  \end{aligned}
\]
The cancellation of terms proportional to \(f^{\mu\nu}\lambda\) requires
that \(c_\lambda=1/2\), while the cancellation of terms proportional to
\(D\lambda\) requires that \(c_D=-c_\lambda\), so the supersymmetric
Lagrangian density takes the form
\begin{equation}
  \mathcal{L}_{\mathrm{gauge}}
  =-\frac14f_{\mu\nu}f^{\mu\nu}
  +\ii\bar\lambda\bar\sigma^\mu\partial_\mu\lambda
  +\frac12D^2.
  \tag{27.2.6}\label{eq:27.2.6}
\end{equation}
This shows that with \(V^\mu\) canonically normalized, the field
\(\lambda\) related to \(V^\mu\) by the transformation rules
\eqref{eq:27.2.2} and \eqref{eq:27.2.3} is also canonically normalized.

In addition, for Abelian gauge theories there is a superrenormalizable
term, known as a \textit{Fayet--Iliopoulos term}%
~\hyperref[ch27-ref-2]{[2]}:
\begin{equation}
  \mathcal{L}_{\mathrm{FI}}=\xi D,
  \tag{27.2.7}\label{eq:27.2.7}
\end{equation}
with \(\xi\) an arbitrary constant. Its variation under a supersymmetry
transformation is shown by Eq.~\eqref{eq:27.2.4} to be a derivative, so
that it yields another supersymmetric term in the action. As we will see
in Section 27.5, the presence of such a term can provide a mechanism for
the spontaneous breakdown of supersymmetry.

Both for its own interest and as a tool in constructing supersymmetric
interactions involving the fields \(f_{\mu\nu}\), \(\lambda\), and \(D\),
it is interesting to ask what sort of superfield has these as component
fields. Somewhat surprisingly, it turns out to be a \textit{spinor
superfield} \(W_\alpha(x)\), with component fields (in the notation of
Eq.~\eqref{eq:26.2.10}) given by
\begin{equation}
  \begin{aligned}
  W_\alpha\big|_{\theta=\bar\theta=0}
    &=\lambda_\alpha(x),\\
  \frac{\partial W_\alpha}{\partial\theta^\beta}
    \bigg|_{\theta=\bar\theta=0}
    &=-(\sigma^{\mu\nu})_\alpha{}_\beta f_{\mu\nu}(x)
      -\ii\delta_\alpha{}^\beta D(x),\\
  W_\alpha\big|_{\theta^2}
    &=-\ii\bigl(\sigma^\mu\partial_\mu\bar\lambda\bigr)_\alpha .
  \end{aligned}
  \tag{27.2.8}\label{eq:27.2.8}
\end{equation}
The assignment \eqref{eq:27.2.8}, written in the two-component form of
the general expansion \eqref{eq:26.2.10}, can be checked directly from
the transformation laws \eqref{eq:26.2.11}--\eqref{eq:26.2.17}, using
in particular the field-strength and auxiliary-field variations
\eqref{eq:27.2.2} and \eqref{eq:27.2.4}.
Substituting \eqref{eq:27.2.8} into that chiral expansion and using the
two-component reduction of the identity \eqref{eq:26.A.5} gives the full
spinor superfield.

Writing \(x_+^\mu=x^\mu-\ii\theta\sigma^\mu\bar\theta\), the full superfield
takes the form
\begin{equation}
  W_\alpha(x,\theta,\bar\theta)
  =\exp\!\left(-\ii\theta\sigma^\mu\bar\theta\,\partial_\mu\right)
   \left[
    \lambda_\alpha
    -(\sigma^{\mu\nu}\theta)_\alpha f_{\mu\nu}
    -\ii\theta_\alpha D
    -\ii\theta^2(\sigma^\mu\partial_\mu\bar\lambda)_\alpha
   \right](x).
  \tag{27.2.9}\label{eq:27.2.9}
\end{equation}
As we showed in Section 26.3, the undotted field strength and its dotted
conjugate form a chiral pair:
\begin{equation}
  W_\alpha(x,\theta,\bar\theta),
  \qquad
  \bar W_{\dot\alpha}(x,\theta,\bar\theta).
  \tag{27.2.10}\label{eq:27.2.10}
\end{equation}
Their component expansions are
\begin{equation}
  \begin{aligned}
  W_\alpha(x,\theta,\bar\theta)
  ={}&\lambda_\alpha(x_+)
    -(\sigma^{\mu\nu}\theta)_\alpha f_{\mu\nu}(x_+)
    -\ii\theta_\alpha D(x_+)
    -\ii\theta^2(\sigma^\mu\partial_\mu\bar\lambda)_\alpha(x_+).
  \end{aligned}
  \tag{27.2.11}\label{eq:27.2.11}
\end{equation}
\begin{equation}
  \begin{aligned}
  \bar W_{\dot\alpha}(x,\theta,\bar\theta)
  ={}&\bar\lambda_{\dot\alpha}(x_-)
    -(\bar\sigma^{\mu\nu}\bar\theta)_{\dot\alpha}
      f_{\mu\nu}(x_-)
    +\ii\bar\theta_{\dot\alpha}D(x_-)
    +\ii\bar\theta^2
      (\bar\sigma^\mu\partial_\mu\lambda)_{\dot\alpha}(x_-),
  \end{aligned}
  \tag{27.2.12}\label{eq:27.2.12}
\end{equation}
with \(x_\pm^\mu\) given by Eq.~\eqref{eq:26.3.23}.

As we saw in Section 26.3, we can construct suitable Lagrangian densities
from the \(\mathcal{F}\)-term of any scalar function of a left-chiral
superfield, plus its Hermitian adjoint. The simplest scalar function of the
left-chiral superfield \eqref{eq:27.2.11} is
\(W^\alpha W_\alpha\).
To calculate the \(\mathcal{F}\)-term, we note that, when expressed as a
function of \(\theta\) and \(x_+\), its term of second order in \(\theta\)
is
\[
  \begin{aligned}
  -[W^\alpha W_\alpha]_{\theta^2}
  =\theta^2\left[
    2\ii\bar\lambda\bar\sigma^\mu\partial_\mu\lambda
    -\frac12f_{\mu\nu}f^{\mu\nu}
    +\frac{\ii}{4}\epsilon^{\mu\nu\rho\sigma}
       f_{\mu\nu}f_{\rho\sigma}
    +D^2\right].
  \end{aligned}
\]
(The argument of the fields here can be taken as \(x^\mu\) rather than
\(x_+^\mu\) because the difference would yield terms that contain at least
three factors of \(\theta\), which therefore vanish.) This follows directly
from
\[
  \theta_\alpha\theta_\beta
  =\frac12\epsilon_{\alpha\beta}\theta^2,
  \qquad
  \theta^\alpha\theta^\beta
  =-\frac12\epsilon^{\alpha\beta}\theta^2,
\]
together with the definition of \(\sigma^{\mu\nu}\). The
\(\mathcal{F}\)-term is the coefficient of \(\theta^2\), so
\begin{equation}
  \begin{aligned}
  -[W^\alpha W_\alpha]_{\mathcal{F}}
  ={}&2\ii\bar\lambda\bar\sigma^\mu\partial_\mu\lambda
    -\frac12f_{\mu\nu}f^{\mu\nu}
  \\
  &+\frac{\ii}{4}\epsilon^{\mu\nu\rho\sigma}
    f_{\mu\nu}f_{\rho\sigma}+D^2.
  \end{aligned}
  \tag{27.2.13}\label{eq:27.2.13}
\end{equation}
The kinetic term is real up to a total derivative, by the two-component
form of Eq.~\eqref{eq:26.A.21}, so the real
part of Eq.~\eqref{eq:27.2.13} yields the Lagrangian
\eqref{eq:27.2.6} for the gauge and gaugino fields
\begin{equation}
  -\frac12\operatorname{Re}\left[
    W^\alpha W_\alpha
  \right]_{\mathcal{F}}
  =\ii\bar\lambda\bar\sigma^\mu\partial_\mu\lambda
  -\frac14f_{\mu\nu}f^{\mu\nu}
  +\frac12D^2.
  \tag{27.2.14}\label{eq:27.2.14}
\end{equation}
The physical significance of the imaginary part will be discussed in a
more general context in the next section.

There is another way to derive the form of the spinor superfield, which
will turn out to provide a more convenient way of deriving the components
of the gauge superfield in non-Abelian gauge theories. A tedious but
straightforward calculation shows that the gauge-invariant superfield
\eqref{eq:27.2.9} may be expressed in terms of the gauge superfield
\eqref{eq:27.1.16} as
\begin{equation}
  W_\alpha(x,\theta,\bar\theta)
  =\frac{\ii}{4}\bar D^2D_\alpha V(x,\theta,\bar\theta),
  \tag{27.2.15}\label{eq:27.2.15}
\end{equation}
where \(D_\alpha\) and \(\bar D_{\dot\alpha}\) are the superderivatives
introduced in Eq.~\eqref{eq:26.2.26}:
\[
  D_\alpha=-\frac{\partial}{\partial\theta^\alpha}
    +\ii\sigma^\mu_{\alpha\dot\alpha}\bar\theta^{\dot\alpha}\partial_\mu,
  \qquad
  \bar D_{\dot\alpha}
    =\frac{\partial}{\partial\bar\theta^{\dot\alpha}}
      -\ii\theta^\alpha\sigma^\mu_{\alpha\dot\alpha}\partial_\mu,
\]
and \(\bar D^2\equiv\bar D_{\dot\alpha}\bar D^{\dot\alpha}\).
This result might have been obtained (aside from the normalization factor)
by noting that the function \eqref{eq:27.2.15} has the desired property of
being a gauge-invariant chiral spinor superfield. First, note that
Eq.~\eqref{eq:27.2.15} is a superfield, because it is formed by acting on
the superfield \(V\) with superderivatives. Also, it follows from the
anticommutation of the superderivatives that any product of three
undotted derivatives, or of three dotted derivatives, vanishes:
\begin{equation}
  D^2D_\alpha=0,
  \qquad
  \bar D^2\bar D_{\dot\alpha}=0.
  \tag{27.2.16}\label{eq:27.2.16}
\end{equation}
The undotted superfield in Eq.~\eqref{eq:27.2.15} and its dotted conjugate
are therefore chiral, with
\begin{equation}
  \begin{aligned}
  W_\alpha(x,\theta,\bar\theta)
  &=\frac{\ii}{4}\bar D^2D_\alpha V(x,\theta,\bar\theta),
  \\
  \bar W_{\dot\alpha}(x,\theta,\bar\theta)
  &=\frac{\ii}{4}D^2\bar D_{\dot\alpha}V(x,\theta,\bar\theta).
  \end{aligned}
  \tag{27.2.17}\label{eq:27.2.17}
\end{equation}
Finally, we can show that \eqref{eq:27.2.15} is invariant under the
generalized gauge transformation \eqref{eq:27.1.13}, which for a single
Abelian gauge superfield is simply
\begin{equation}
  V(x,\theta,\bar\theta)\longrightarrow
  V(x,\theta,\bar\theta)+\frac{\ii}{2}
  \left[\Omega(x,\theta,\bar\theta)
    -\Omega^*(x,\theta,\bar\theta)\right],
  \tag{27.2.18}\label{eq:27.2.18}
\end{equation}
where \(\Omega\) is an arbitrary left-chiral superfield. Since
\(D_\alpha\Omega^*=0\), the change in \(W_\alpha\) is proportional to
\(\bar D^2D_\alpha\Omega\). But \(\bar D_{\dot\alpha}\Omega=0\) and
\[
  [\bar D^2,D_\alpha]
  =-4\ii\sigma^\mu_{\alpha\dot\alpha}
    \partial_\mu\bar D^{\dot\alpha},
\]
so the change in \(W_\alpha\) vanishes. A similar argument shows that
\(\bar W_{\dot\alpha}\) is also gauge-invariant. (The work of checking
Eq.~\eqref{eq:27.2.15} is greatly reduced by using this gauge-invariance
property to put \(V(x,\theta)\) in Wess--Zumino gauge.)

The chiral superfields \eqref{eq:27.2.11} and \eqref{eq:27.2.12} are
evidently not of the most general form for undotted and dotted chiral
superfields. To put the constraints satisfied by these superfields in a
manifestly supersymmetric form, we note by using the anticommutation
relation \eqref{eq:26.2.30} that
\begin{equation}
  D^\alpha\bar D^2D_\alpha
  =\bar D_{\dot\alpha}D^2\bar D^{\dot\alpha}.
  \tag{27.2.19}\label{eq:27.2.19}
\end{equation}
From Eq.~\eqref{eq:27.2.17} it follows then that \(W_\alpha\) and
\(\bar W_{\dot\alpha}\) are
related by the constraint
\begin{equation}
  D^\alpha W_\alpha
  =\bar D_{\dot\alpha}\bar W^{\dot\alpha}.
  \tag{27.2.20}\label{eq:27.2.20}
\end{equation}
It is straightforward to show that the most general chiral spinor
superfields satisfying Eq.~\eqref{eq:27.2.20} are of the form
\eqref{eq:27.2.11} and \eqref{eq:27.2.12}, with \(f_{\mu\nu}\) constrained
to satisfy the `Bianchi' identities
\(\epsilon^{\mu\nu\rho\sigma}\partial_\rho f_{\mu\nu}=0\).
